\section*{ID6830: Tidal-Tensor/Filament Alignment Test: λ\_2>λ\_rm Th}
\paragraph{Metadata.}
PiID: 2408514006660 | RTEID: RTE:P26522.4068:T6.0017 | UUID: 753b1cc3-cce8-508f-8147-b99df6ae16e8

\paragraph{Canonical Statement.}
\#\#\# B. Tidal-tensor/filament alignment test 1. From initial potential estimate (linear-regime reconstruction or simulation), compute \textbackslash{}(T\_\{ij\}(\textbackslash{}mathbf\{q\})=\textbackslash{}partial\_i\textbackslash{}partial\_j\textbackslash{}Phi\textbackslash{}). 2. For each spatial cell find eigenvectors \textbackslash{}(\textbackslash{}mathbf\{e\}\_1,\textbackslash{}mathbf\{e\}\_2,\textbackslash{}mathbf\{e\}\_3\textbackslash{}) and eigenvalues \textbackslash{}(\textbackslash{}lambda\_1\textbackslash{}ge\textbackslash{}lambda\_2\textbackslash{}ge\textbackslash{}lambda\_3\textbackslash{}). 3. Flag filament candidate cells where \textbackslash{}(\textbackslash{}lambda\_1>\textbackslash{}lambda\_\{\textbackslash{}rm th\}\textbackslash{}) and \textbackslash{}(\textbackslash{}lambda\_2>\textbackslash{}lambda\_\{\textbackslash{}rm th\}\textbackslash{}) but \textbackslash{}(\textbackslash{}lambda\_3<\textbackslash{}lambda\_\{\textbackslash{}rm th\}\textbackslash{}). 4. For these cells, compute the alignment angle between \textbackslash{}(\textbackslash{}mathbf\{e\}\_3\textbackslash{}) (the filament direction) and the direction(s) where the angular spectral density (from SBT) is maximal. High alignment → supports the modal-origin hypothesis.

\paragraph{Canonical Equation.}
\[
\lambda_2>\lambda_{\rm th}
\]

\paragraph{Dictionary.}
Tidal-Tensor/Filament Alignment Test: λ\_2>λ\_rm Th — λ\_2>λ\_rm th

\paragraph{Encyclopedia.}
Tidal-Tensor/Filament Alignment Test: λ\_2>λ\_rm Th is an inequality / bound, written λ\_2>λ\_rm th. It states an inequality or bound that constrains the admissible range of the quantity. It is built from h. Within the Trawin Zero Point registry it serves as one element of the framework's mathematical narrative.
